数值方法默认解就在那里，问题只是怎样算准。这一单元退后一步，问三个更靠前的问题：变化规律怎样翻译成方程，什么样的方程能被完全解出，以及解到底存不存在、唯不唯一。

整个单元的落差只有一句话：微分方程交给你的是每一点上的变化率，你要的是穿过这些变化率的那条完整轨迹。

\begin{center}
\begin{tikzpicture}[>=stealth]
  \foreach \tt in {0.3,0.9,1.5,2.1}{
    \foreach \xx/\ss in {0.3/0.55, 0.9/0.35, 1.5/0.1, 2.1/-0.2}{
      \draw[thin,gray,->] ({\tt-0.16},{\xx-0.16*\ss}) -- ({\tt+0.16},{\xx+0.16*\ss});
    }
  }
  \draw[->] (2.7,1.2) -- (3.5,1.2);
  \node[above] at (3.1,1.25) {\small 求解};
  \begin{scope}[shift={(4.0,0)}]
    \foreach \tt in {0.3,0.9,1.5,2.1}{
      \foreach \xx/\ss in {0.3/0.55, 0.9/0.35, 1.5/0.1, 2.1/-0.2}{
        \draw[thin,gray,->] ({\tt-0.16},{\xx-0.16*\ss}) -- ({\tt+0.16},{\xx+0.16*\ss});
      }
    }
    \draw[very thick] plot[smooth] coordinates
      {(0.1,0.15) (0.6,0.45) (1.1,0.85) (1.6,1.25) (2.1,1.58) (2.5,1.75)};
  \end{scope}
\end{tikzpicture}
\end{center}

\emph{左边是方程给的局部规则，右边是要求的整体轨迹；本单元的四节都在这条箭头上。}

第一节从增长、衰减与振动三个经典模型出发，说明微分方程是变化规律的精确化，并用方向场给出不解方程也能看见解的几何眼光——图上灰色小箭头正是那样一张方向场。第二节把能彻底解出的那一类做成样板：\(x'=Ax\) 的解写成 \(e^{tA}x_0\)，可对角化时就是特征模态的叠加，长期行为整个由特征值读出。第三节走到地基处，用有限时间爆破和无穷多解两个灾难逼出 Picard--Lindelöf 定理，逐条交代每个条件挡住哪种失败。第四节处理解不出来的情形：无量纲化找出小参数，正则摄动逐阶修正，快慢分离让快变量退化为瞬时平衡，而小参数一旦压在最高阶导数上就会冒出边界层。

主线因此是：先用可解的例子建立直觉，再把线性系统做成样板，用地基性的定理划清``问题提对了''的边界，最后在无法精确求解时借尺度差异构造带适用范围的近似。后面讨论稳定性与相图时，默认的存在唯一性前提都出自本单元；控制模型里的快慢分解也会再次用到渐近观点。

\subsection{怎么读这一章}

首次阅读抓住前三节的判断链就够：变化规律怎样变成初值问题，线性系统怎样由谱读出行为，连续性与 Lipschitz 条件分别保证什么。第四节适合第二遍读，重点不在记住摄动展开的写法，而在养成两个习惯——先无量纲化，再检查小参数是否乘在最高阶导数前。存在唯一性定理的完整证明可以后置，但那两个反例和条件对账不能跳过。

\subsection{本章篇目}

以下篇目按概念依赖排列；若从具体问题进入，也可以先打开最接近当前困惑的一篇，再沿篇内链接回补。

\begin{itemize}
\tightlist
\item \hyperref[sec:15-Differential-Equations-and-Dynamical-Systems/01-ODE-Theory/01]{``从增长、衰减与振动说起''}；
\item \hyperref[sec:15-Differential-Equations-and-Dynamical-Systems/01-ODE-Theory/02]{``线性系统与矩阵指数''}；
\item \hyperref[sec:15-Differential-Equations-and-Dynamical-Systems/01-ODE-Theory/03]{``存在唯一性为什么需要 Lipschitz''}；
\item \hyperref[sec:15-Differential-Equations-and-Dynamical-Systems/01-ODE-Theory/04]{``尺度分离与渐近近似''}。
\end{itemize}

读完这四节，手里握着的是``解存在且唯一''这块地基，以及一套只对线性系统完全奏效的求解术。下一单元换一副眼光：不去求解，直接在状态空间里看轨迹的长期走向。
